\section{模型检验、稳健性与局限性}

\subsection{分层验证设计}

本文依据各模块的信息结构，按“数据层—代理层—机理层—决策层”分别设置验证指标。

\begin{table}[H]
\centering
\caption{验证方法与结果}
\label{tab:validation}
\small
\begin{tabularx}{0.97\textwidth}{C{2.3cm}L L C{2.4cm}}
\toprule
模块 & 验证方法 & 评价量 & 结果\\
\midrule
数据质量 & 时间网格与字段一致性检验 & 行数、缺失、边界计数 & 全部通过\\
删失估计 & $n=50,000$ 合成正态数据，$\mu=50.15,\sigma=0.30$ & 参数恢复误差 & $|\Delta\mu|,|\Delta\sigma|<0.03$\\
能耗代理 & 逐日留一交叉验证 & 平均 $R^2$/RMSE & 0.9866 / 5.96 kW\\
机理方向 & 单元测试提高电压、延长周期 & 浓度变化符号 & 均符合预期\\
机会约束 & 独立 10,000 情景复核 & 95\%分位、达标概率 & $10\mgNm$ 八工况全可行\\
分层优先级 & 约束容差与 12 组独立情景 & R2/R8 优先类别 & 电压优先/周期优先均为 12/12\\
问题四稳定性 & 110\% 宽边界下 R8 的 12 组独立重采样 & 增幅经验 95\% 范围 & $[4.53\%,6.13\%]$\\
\bottomrule
\end{tabularx}
\end{table}

能耗验证采用整日分组，以隔离分钟级自相关造成的信息泄漏；排放验证采用删失似然、物理方向检查和独立机会约束复核。11 项自动化检验覆盖数据不变量、合成删失恢复、能耗精度、物理方向、边界反向差分、活跃约束分类和两级限值可行性逻辑。在 $10^{-6}$--$10^{-3}$ 的数值容差范围内，R1--R3 始终归入电压优先层，R4--R8 始终归入周期优先层；各周期变量的 $G_{T_i}$ 均为正，电压变量对 $I_M$ 的直接松弛收益为零。

\subsection{模型性质}

\begin{enumerate}
  \item 在建模前分析删失特征与闭环控制偏差，减少输出信息不足造成的参数误判；
  \item 公式中每个控制方向有明确工程含义，能输出场间归因和控制优先级；
  \item 工况识别、能耗拟合、优化和独立验证均保持时间结构，计算过程可以重复；
  \item 对无法满足 $5\mgNm$ 约束的工况单独报告，并进一步分析设备能力扩展情景；
  \item 机会约束同时包含工况波动、模型残差和参数不确定性，比均值约束更适合环保风险控制。
\end{enumerate}

\subsection{局限与改进方向}

\begin{enumerate}
  \item \textbf{排放量程。} 54.75\% 的有效读数封顶且目标限值在观测域外，是最主要局限。应配置 0--20 或 0--30 $\mgNm$ 的高精度在线仪表，并保留原量程作冗余。
  \item \textbf{电气变量。} 缺少二次电流、火花率和反电晕状态，而这些量是电除尘运行检查的重要信息\cite{szabo1981inspection}。本文建立的电压平方模型仅适用于历史工作区间；加入上述变量后，可进一步建立伏安特性与火花率约束。
  \item \textbf{粉尘性质。} 粉尘比电阻、粒径、含湿量和化学成分未测，温度效应只能用先验表示；后续可用分层贝叶斯模型按原料批次更新。
  \item \textbf{振打相位。} 由于缺少实际触发时刻，本文采用相位积分描述平均效应并给出错峰规则；接入 PLC 事件记录后，可利用脉冲响应模型，或采用能够描述事件相互激发关系的 Hawkes 模型识别再飞扬峰值。
  \item \textbf{短期数据。} 七天数据覆盖的季节与设备状态有限。部署时应采用滚动校准、漂移检测和策略回退机制。
\end{enumerate}

\subsection{工程应用前的最小标定试验}

在满足环保与设备约束的前提下，可选择 2--3 个稳定工况进行小幅阶跃试验：每次仅改变一场电压 1\%--2\%，并分别记录各场振打事件；同步采集低量程出口浓度、二次电流、火花率和灰斗料位。利用这些数据重新估计 $D_0,\bm w,\kappa_M,\delta$，并进行不少于一周的离线跟踪验证。仅当独立样本的 95\% 覆盖率满足要求且策略切换频率符合设备约束时，方可进入闭环运行。
